\documentclass[11pt]{article}

\usepackage[letterpaper,margin=1in]{geometry}
\usepackage{amsmath,amssymb}
\usepackage{booktabs}
\usepackage{longtable}
\usepackage{array}
\usepackage{xcolor}
\usepackage{listings}
\usepackage{url}
\usepackage{hyperref}
\usepackage{titlesec}
\usepackage{parskip}
\usepackage{fancyhdr}
\usepackage{enumitem}

\hypersetup{
  colorlinks=true,
  linkcolor=black,
  urlcolor=blue!50!black,
  citecolor=black,
  pdftitle={KV Cache Engine --- Software Overview},
  pdfauthor={LonghornSilicon},
  pdfsubject={Compiler Co-Design Orientation}
}

\titleformat{\section}{\normalfont\Large\bfseries}{\thesection.}{0.6em}{}
\titleformat{\subsection}{\normalfont\large\bfseries}{\thesubsection}{0.6em}{}

\pagestyle{fancy}
\fancyhf{}
\lhead{\small KV Cache Engine --- Software Overview}
\rhead{\small v0.2}
\cfoot{\thepage}
\renewcommand{\headrulewidth}{0.4pt}

\definecolor{codebg}{gray}{0.97}
\definecolor{codekey}{rgb}{0.13,0.29,0.53}
\definecolor{codecom}{rgb}{0.25,0.5,0.25}
\lstset{
  basicstyle=\ttfamily\small,
  backgroundcolor=\color{codebg},
  frame=single,
  framerule=0.4pt,
  rulecolor=\color{black!30},
  xleftmargin=0.5em,
  xrightmargin=0.5em,
  aboveskip=0.6em,
  belowskip=0.6em,
  showstringspaces=false,
  breaklines=true,
  keywordstyle=\color{codekey}\bfseries,
  commentstyle=\color{codecom}\itshape,
  numbers=none
}

\setlist{itemsep=1pt,topsep=2pt}

\begin{document}

%==============================================================================
\thispagestyle{empty}
\begin{center}
\vspace*{1.2in}
{\Huge\bfseries KV Cache Engine}\\[0.4em]
{\LARGE\bfseries Software Overview}\\[1.2em]

{\large \textsc{LonghornSilicon --- Compiler Co-Design Orientation}}\\[0.4em]
{\large Block 2 of 4 --- KV Cache Compression / Decompression}\\[2em]

\rule{0.7\textwidth}{0.4pt}\\[1em]

\begin{tabular}{rl}
\textbf{Version}:    & 0.2 (ChannelQuant codec) \\
\textbf{Date}:       & 2026-07-03 \\
\textbf{Audience}:   & Compiler team, runtime engineers \\
\end{tabular}\\[1em]
\rule{0.7\textwidth}{0.4pt}
\vfill
\end{center}
\clearpage

%==============================================================================
\tableofcontents
\clearpage

%==============================================================================
\section{Purpose}

This document orients the compiler team on the KV cache engine block.
It covers: what the block does, how to use the reference models, what
the compiler should emit, and what integration phases look like.

For the formal interface contract, see the ISA specification
(\path{docs/isa/kv_cache_engine_isa.pdf}). For the API reference, see
\path{docs/reference_model_api.pdf}.

%==============================================================================
\section{What the Block Does}

The KV cache engine sits between the attention compute unit (block~1)
and the memory hierarchy controller (block~4). It compresses key and
value vectors before writing to on-chip SRAM, and decompresses them
when the attention unit needs them for dot-product computation.

\begin{itemize}
\item \textbf{On write}: incoming FP16 K/V vectors are compressed via
      \textbf{ChannelQuant} (per-channel INT4 keys grouped over $G$ tokens
      + optional FP16 outlier channels; per-token INT4 values) and stored in SRAM.
\item \textbf{On read}: compressed entries are decompressed back to
      \textbf{fp32} coordinates and streamed to the attention unit via AXI-Stream.
\item \textbf{Eviction}: when SRAM is full, the block signals the memory
      hierarchy controller to spill entries to DRAM.
\end{itemize}

\subsection{Compression Pipeline}

\begin{enumerate}
\item \textbf{Key group buffer} --- accumulate $G$ key tokens
      (\texttt{residual\_buffer}).
\item \textbf{Per-channel scale} --- per-channel max over the group $\rightarrow$
      $D$ frozen FP16 scales (\texttt{amax\_unit} + \texttt{scale\_bank}); values
      use a per-token FP16 scale.
\item \textbf{Quantize} --- keep channels $\rightarrow$ INT4; the top-$k$ outlier
      channels (CQ-4+) are held FP16 from a static ROM mask. Values $\rightarrow$
      INT4 (INT8 for the CQ-8 tier).
\item \textbf{Store} --- a unified per-channel SRAM record (keep $\to\{$scale,
      INT4$\}$; outlier $\to\{$fp16, code $+1\}$).
\end{enumerate}

\noindent Decompression is per channel: $\hat{x}_c = \text{code}_c \cdot
\text{scale}_c$ in fp32; outlier channels replay their stored FP16 directly. No
rotation, no Johnson--Lindenstrauss, no Lloyd--Max codebook (those were the
retired TurboQuant+ codec).

%==============================================================================
\section{Reference Models}

\subsection{Directory Layout}

\begin{lstlisting}[language=bash]
sw/
+-- README.md                     <- orientation (this doc's markdown form)
+-- reference_model/
    +-- README.md                 <- API reference + build instructions
    +-- Makefile
    +-- channelquant_ref.hpp/.cpp <- ChannelQuant C++ reference (shipped codec)
    +-- test_channelquant_ref.cpp <- 9 golden-vector tests
    +-- kv_cache_engine_ref.*     <- retired TurboQuant+ reference (archival only)
# Python algorithm reference + golden vectors: ../channelquant/reference/
\end{lstlisting}

\subsection{Which Models Are Available}

\begin{table}[h]
\centering
\small
\begin{tabular}{@{}l l l l@{}}
\toprule
Codec & Status & C++ API & Python \\
\midrule
ChannelQuant (shipped) & Bit-exact vs RTL + golden & \texttt{lhsi::cq::compress\_*} & \texttt{../channelquant} ref \\
TurboQuant+ (retired) & archival regression only & \texttt{lhsi::KVCacheEngine} & --- \\
\bottomrule
\end{tabular}
\end{table}

\subsection{Build Everything}

\begin{lstlisting}[language=bash]
cd sw/reference_model
make test-cq      # ChannelQuant C++ tests vs the 9 golden vectors
make test-all     # legacy TQ + ChannelQuant + Python tests
make shared       # libkv_cache_engine_ref.so
make static       # libkv_cache_engine_ref.a
\end{lstlisting}

\noindent Requires Python 3.10+, NumPy, and a C++17 compiler.

%==============================================================================
\section{What the Compiler Should Emit}

A compiler backend targeting the KV cache engine emits sequences of
the operations defined in the ISA spec (\S4 of
\path{docs/isa/kv_cache_engine_isa.pdf}):

\begin{table}[h]
\centering
\small
\begin{tabular}{@{}l l p{6cm}@{}}
\toprule
Operation & C API & What it does \\
\midrule
\texttt{KV\_QUERY} & \texttt{lhsi\_kv\_info()} & Read hardware config \\
\texttt{KV\_RESET} & \texttt{lhsi\_kv\_reset()} & Clear SRAM, reset FSM \\
\texttt{KV\_COMPRESS\_KEY} & \texttt{lhsi\_kv\_compress\_key()} & Compress + store key \\
\texttt{KV\_COMPRESS\_VALUE} & \texttt{lhsi\_kv\_compress\_value()} & Compress + store value \\
\texttt{KV\_DECOMPRESS\_KEY} & \texttt{lhsi\_kv\_decompress\_key()} & Read + decompress key \\
\texttt{KV\_DECOMPRESS\_VALUE} & \texttt{lhsi\_kv\_decompress\_value()} & Read + decompress value \\
\bottomrule
\end{tabular}
\end{table}

\subsection{Worked Example: Attention with KV Cache}

The shipped reference is head-level ChannelQuant (frozen in \path{../channelquant}):
compress a whole KV head $[T, D]$, then reconstruct for attention.

\begin{lstlisting}[language=Python]
import channelquant_ref as cq   # ../channelquant/reference

# Store phase: per-channel INT4 keys (grouped G), per-token INT4 values.
K_hat = cq.fq_per_channel(K, bits=4, G=128)          # or fq_per_channel_outlier for CQ-4+
V_hat = cq.fq_per_token(V, bits=4)

# Attention phase: use the reconstructed K_hat / V_hat.
scores = query @ K_hat.transpose(-1, -2) / sqrt(D)   # Q * K^T
# Route each S*V tile through the precision controller (block 1); see below.
\end{lstlisting}

%==============================================================================
\section{Integration Phases}

\begin{table}[h]
\centering
\small
\renewcommand{\arraystretch}{1.2}
\begin{tabular}{@{}c p{3cm} p{4cm} p{4.5cm}@{}}
\toprule
Phase & Timeline & Compiler targets & We provide \\
\midrule
0 & now & Python + C++ reference models & Models + ISA + tests \\
1 & when FPGA arrives & AXI-Lite + AXI-Stream on Zynq & Vivado bitstream + driver \\
2 & all 4 blocks done & Multi-block FPGA project & Integrated attention pipeline \\
3 & post-tape-out & PCIe-attached chip & TSMC 16FFC silicon + driver \\
\bottomrule
\end{tabular}
\end{table}

\noindent The ISA contract (\path{docs/isa/kv_cache_engine_isa.pdf}) is
stable across all four phases. Only the runtime implementation changes ---
Python function call, AXI register write, or PCIe MMIO.

%==============================================================================
\section{Interaction with Block~1 (Precision Controller)}

The KV cache engine and precision controller work together during
attention inference:

\begin{enumerate}
\item KV cache engine decompresses K/V vectors from SRAM
\item The attention unit computes $Q \cdot K^T$ scores
\item The precision controller decides INT8 vs FP16 per tile
\item The MAC array computes $\mathrm{softmax}(S) \cdot V$ at the chosen precision
\end{enumerate}

\noindent A compiler backend that targets both blocks emits interleaved
\texttt{KV\_DECOMPRESS\_*} and \texttt{PC\_PUSH\_TILE} operations. The
reference models for both blocks can be composed in the same process:

\begin{lstlisting}[language=Python]
import channelquant_ref as cq                        # block 2 (KVE)
from precision_controller_ref import PrecisionController  # block 1 (ACU)

pc = PrecisionController()

K_hat = cq.fq_per_channel(K, bits=4, G=128)          # decompressed keys
scores = query @ K_hat.transpose(-1, -2)             # Q * K^T
decision = pc.process_tile(scores)                   # True = FP16, False = INT8
\end{lstlisting}

%==============================================================================
\section{Open Questions for the Compiler Team}

\begin{enumerate}
\item \textbf{Granularity}: does the compiler emit individual compress/decompress
      ops, or a fused \texttt{kv\_attention(Q, K\_cache, V\_cache)} op?
\item \textbf{Batch interface}: should the reference model support batched
      compress/decompress (multiple vectors in one call) for throughput?
\item \textbf{Async model}: should \texttt{KV\_COMPRESS\_KEY} be non-blocking
      with a separate completion query, matching the hardware pipeline?
\item \textbf{Cost model}: are placeholder cycle counts sufficient, or do
      you need post-synthesis numbers for the scheduler?
\end{enumerate}

\vspace{0.5in}
\noindent\rule{\textwidth}{0.4pt}\\
\noindent\small\textit{LonghornSilicon KV Cache Engine --- Software Overview.
This document is open and may be redistributed.}

\end{document}
